\section*{Sets and Indices}

There are no nontrivial index sets in the implemented model. The problem uses two truck types:
\[
\{\text{Type A}, \text{Type B}\}.
\]

\section*{Parameters}

From \texttt{general\_parameters.csv}, let:
\begin{align*}
R_A &:= \text{refrigerated capacity of a Type A truck} 
&& (\texttt{refrigerated\_capacity\_type\_a}) = 20,\\
N_A &:= \text{non-refrigerated capacity of a Type A truck} 
&& (\texttt{non\_refrigerated\_capacity\_type\_a}) = 40,\\
R_B &:= \text{refrigerated capacity of a Type B truck} 
&& (\texttt{refrigerated\_capacity\_type\_b}) = 30,\\
N_B &:= \text{non-refrigerated capacity of a Type B truck} 
&& (\texttt{non\_refrigerated\_capacity\_type\_b}) = 30,\\
D_R &:= \text{required refrigerated cargo volume} 
&& (\texttt{refrigerated\_cargo\_requirement}) = 3000,\\
D_N &:= \text{required non-refrigerated cargo volume} 
&& (\texttt{non\_refrigerated\_cargo\_requirement}) = 4000,\\
c_A &:= \text{rental cost per Type A truck per kilometer} 
&& (\texttt{rental\_cost\_type\_a}) = 30,\\
c_B &:= \text{rental cost per Type B truck per kilometer} 
&& (\texttt{rental\_cost\_type\_b}) = 40.
\end{align*}

\section*{Decision Variables}

\begin{align*}
x &:= \text{number of Type A trucks rented} && (\texttt{x}) \in \mathbb{Z}_{\ge 0},\\
y &:= \text{number of Type B trucks rented} && (\texttt{y}) \in \mathbb{Z}_{\ge 0}.
\end{align*}

\section*{Objective}

Minimize total rental cost per kilometer:
\[
\min \; c_A x + c_B y.
\]

With the instance data, this is
\[
\min \; 30x + 40y.
\]

\section*{Constraints}

Refrigerated cargo requirement:
\[
R_A x + R_B y \ge D_R.
\]

Non-refrigerated cargo requirement:
\[
N_A x + N_B y \ge D_N.
\]

Integrality and nonnegativity:
\[
x, y \in \mathbb{Z}_{\ge 0}.
\]

With the instance data, the capacity constraints are
\begin{align*}
20x + 30y &\ge 3000,\\
40x + 30y &\ge 4000.
\end{align*}

\section*{Notes on Implementation}

The code constructs exactly the above integer linear program in PuLP:
\begin{itemize}
\item \texttt{x} is created as \texttt{LpVariable("TypeA\_trucks", lowBound=0, cat='Integer')}.
\item \texttt{y} is created as \texttt{LpVariable("TypeB\_trucks", lowBound=0, cat='Integer')}.
\item The objective is \texttt{cost\_a * x + cost\_b * y}.
\item The two constraints are added with ``\(\ge\)'' sense exactly as written above.
\end{itemize}

The model is solved using \texttt{pulp.GUROBI\_CMD(msg=0)}. No upper bounds, fixed truck limits, distance scaling, or additional logical constraints are included in the implemented model. The output reports the integer values of \texttt{x} and \texttt{y} and the resulting objective value.
